\documentclass[11pt,a4paper]{article}
\usepackage[T1]{fontenc}
\usepackage{lmodern,amsmath,amssymb,amsthm,mathtools,booktabs,microtype}
\usepackage[margin=27mm]{geometry}
\usepackage[colorlinks=true,linkcolor=blue,citecolor=blue,urlcolor=blue]{hyperref}
\setlength{\emergencystretch}{2em}
\newtheorem{theorem}{Theorem}[section]
\newtheorem{proposition}[theorem]{Proposition}
\newtheorem{lemma}[theorem]{Lemma}
\newtheorem{corollary}[theorem]{Corollary}
\theoremstyle{definition}\newtheorem{definition}[theorem]{Definition}
\newcommand{\Q}{\mathbb Q}\newcommand{\Z}{\mathbb Z}
\DeclareMathOperator{\Tr}{Tr}\DeclareMathOperator{\Nm}{N}
\title{Adaptive levels and varying norm classes for native elliptic detectors}
\author{}\date{16 November 2026}
\begin{document}\maketitle
\begin{abstract}
We give complement levels and varying quadratic norm classes that retain
the negative-odd-valuation implication of the native elliptic detectors.
An exact re-expression leaves the detector output unchanged and makes
every odd valuation of its physical weight positive. For the re-expressed
detector, the level multiplier $P(V,J)/Q(V,J)$ may use polynomials
$P,Q$ with $S$-integral coefficients and an $S$-unit value $Q(0,0)$.
Only that residue configuration is relevant to a negative odd pole. The example
$1/(1-V+J)$ is a unit at every prepared point-algebra character and
can have poles at positive contact valuations. A family of varying
classes has automatic real nonvanishing and the required even local
valuations. All denominators and complement guards are included.
If both actual prescribed-trace norm equations can be realized
simultaneously for every integer input, the existing finite-algebra
encoding gives twenty-one variables for the general source and nine
for either conditional source. That rational realization assertion is
stated explicitly and is not proved; no smaller defining-variable
bound is asserted.
\end{abstract}

\section{Native detectors and arithmetic scope}

Let $K=\Q(\sqrt2)$ and let $S$ be the original finite exceptional
prime set of the chosen native source; in particular $2,3,17\in S$.
We retain the arithmetic premises, sufficient coefficient range,
point-stage guards, local extraction argument, and height argument of
the general twenty-five-variable source or either conditional
thirteen-variable source in~\cite{Unit,Terminal,Complement}.

For one native block write
\begin{equation}\label{eq:native}
 W=\prod_{k=1}^{r}(1+w h_k^2),\qquad
 V=\Tr(W^{-1}),\qquad J=\Nm(W^{-1}),\qquad b=\frac{4w}{V^2}.
\end{equation}
Here $w$ is a nonzero rational scalar depending only on the retained
base coordinates, and $h_k$ are the native field-valued rational
functions. All divisions in~\eqref{eq:native} are guarded.
For the general source $w=\rho R$; for a conditional source
$w=F/(4\delta^8)$, absorbing the fixed guarded square in its
physical input. The factorization~\eqref{eq:native} holds for these
sources as in~\cite{Bounded}.

Let $\mathcal B$ be the finite point algebra at an actual rational
point-stage specialization. It is split by $K$, so every geometric
character of $\mathcal B$ takes values in $K$~\cite[Lemma 1.1]{Complement}.
Trace and norm in~\eqref{eq:native} are the coefficient operations
in $\mathcal B\otimes_\Q K$. At the chosen rational character they
are the usual field trace and norm. At another character the two
slots need not be conjugates; their sum and product still give the
evaluations of $V$ and $J$.

An actual tuple is \emph{prepared} when it has the source's
point-stage guards and prepared detector inputs. In particular the
chosen rational inputs have square reduced denominator, lie in
$(0,17)$, are squares at $2,3,17$, and have the required positive
valuations at those primes. The native scalar $w$ is then positive
and has even $17$-adic valuation. At every geometric character,
every real value of a factor $1+w h_k^2$ is at least one. Thus
both real slots of $W^{-1}$ belong to $(0,1]$.

Every auxiliary function introduced below is a fixed rational
expression in the retained native data. It is not an additional free
rational parameter. A finite list of such expressions can be used
by the finite-union encoding of~\cite{Finite}; no infinite choice of
coefficients or class representatives is implicit.

\section{Local residues and trace-dependent levels}

\begin{lemma}\label{lem:residues}
Let $p\notin S$ and suppose $v_p(w)$ is odd. At an actual rational
character, $V,J$ are $p$-integral and their residue pair is one of
\[
                         (0,0),\qquad(1,0),\qquad(2,1).
\]
At an inert prime the middle pair does not occur.
\end{lemma}
\begin{proof}
A zero $h_k$ contributes the factor one. For every other factor,
at each place of $K$ above $p$, the summands in $1+w h_k^2$
have valuations zero and an odd integer. Their minimum is unique.
The reciprocal of the factor is therefore integral, reducing to
one if the second summand has positive valuation and to zero if it
has negative valuation. Each local slot of $W^{-1}$ consequently
reduces to zero or one. Their sum and product give the displayed
pairs. At an inert prime the two slots are conjugate and have the
same residue in this two-element set.
\end{proof}

\begin{proposition}\label{prop:tracelevels}
Let $P,Q\in\Z[S^{-1}][v,j]$ have nonnegative coefficients and
positive constant coefficients. Suppose each value of $P$ and $Q$
at $(0,0),(1,0),(2,1)$ is an $S$-unit. For a fixed positive integer
$M$, put
\[
                     m=M\frac{P(V,J)}{Q(V,J)},\qquad a=17m^2.
\]
At every $p\notin S$ with $v_p(w)$ odd,
$v_p(m)=v_p(M)\ge0$. Moreover $P(V,J)$, $Q(V,J)$, and $m$
are units at every prepared point-algebra specialization.
\end{proposition}
\begin{proof}
The local assertion follows from Lemma~\ref{lem:residues} and the
six assumed unit values. At every prepared real character, $V,J$
are positive. The coefficient and constant-term hypotheses make
$P(V,J)$ and $Q(V,J)$ strictly positive there. They are therefore
nonzero at every geometric character and hence units in the finite
reduced algebra.
\end{proof}

For example, the multiplier
\begin{equation}\label{eq:boundedexample}
                         m=M\frac{1+V}{1+J}
\end{equation}
has residue ratios $1,2,3/2$. These are units outside $S$.
If $u_+,u_-\in(0,1]$ are the two reciprocal slots, then
$V=u_++u_-$, $J=u_+u_-$, and
\[
 V>J,\qquad 2(1+J)-(1+V)
       =u_+u_-+(1-u_+)(1-u_-)>0.
\]
Consequently $M<m<2M$ at every prepared real character.

\section{A pole-preserving re-expression}

Adjoin the guard $1+w\ne0$ and define
\begin{equation}\label{eq:normalization}
 w^*=\frac{w}{(1+w)^2},\qquad
 h_0^*=1+w,\qquad h_k^*=(1+w)h_k\quad(1\le k\le r).
\end{equation}
This guard holds at every prepared tuple because $w>0$.

\begin{lemma}\label{lem:exact}
The product and reciprocal trace and norm satisfy
\begin{equation}\label{eq:staridentities}
 \begin{gathered}
 W^*:=\prod_{k=0}^{r}(1+w^*(h_k^*)^2)=(1+w)W,\\
 V^*:=\Tr((W^*)^{-1})=\frac{V}{1+w},\qquad
 J^*:=\Nm((W^*)^{-1})=\frac{J}{(1+w)^2},\\
                         \frac{4w^*}{(V^*)^2}=b.
 \end{gathered}
\end{equation}
Every odd valuation of $w^*$ is positive.
\end{lemma}
\begin{proof}
For $k\ge1$, $w^*(h_k^*)^2=w h_k^2$, while the factor with
index zero is $1+w$. Rational scalar multiplication has the displayed
effects on reciprocal trace and norm, proving all identities.

If $v_p(w)>0$, then $v_p(w^*)=v_p(w)$; if $v_p(w)<0$, then
$v_p(w^*)=-v_p(w)$; and if $v_p(w)=0$, then
$v_p(w^*)=-2v_p(1+w)$ is even. Hence every odd value is positive.
\end{proof}

\begin{theorem}[Adaptive trace levels]\label{thm:adaptive}
Let $P,Q\in\Z[S^{-1}][v,j]$, with $Q(0,0)$ an $S$-unit. Set
\begin{equation}\label{eq:adaptive}
                     m=\frac{P(V^*,J^*)}{Q(V^*,J^*)},
                     \qquad a=17m^2,
\end{equation}
and retain the nonzero guards on the numerator and denominator.
If $p\notin S$ and $v_p(b)$ is negative and odd, then
\[
                    v_p(V^*)>0,\qquad v_p(J^*)>0,
                    \qquad v_p(m)\ge0.
\]
In particular, if a rational class value $d\ne0$ has even valuation
at every such prime, rational solubility of
\begin{equation}\label{eq:complement}
                         q^2=d\,a(a-b)
\end{equation}
excludes all negative odd valuations of $b$ outside $S$.
No condition on $P$ or $Q$ at $(1,0)$ or $(2,1)$ is needed.
\end{theorem}
\begin{proof}
The identity $b=4w^*/(V^*)^2$ shows that
$\nu=v_p(w^*)$ is odd, and Lemma~\ref{lem:exact} makes it positive.
The negative valuation of $b$ gives
\[
                         v_p(V^*)>\nu/2>0.
\]
Apply Lemma~\ref{lem:residues} to $W^*$. Its two local reciprocal
slots reduce to zero or one. Their sum is zero modulo $p$, and
$p\ne2$, so both reduce to zero. Their product $J^*$ also reduces
to zero. Thus $P(V^*,J^*)$ is integral, while
$Q(V^*,J^*)\equiv Q(0,0)\not\equiv0\pmod p$.
This proves $v_p(m)\ge0$.

The valuation of $a=17m^2$ is now even and nonnegative, strictly
greater than the negative odd valuation of $b$. Thus $a-b$ has
the valuation of $b$, and the right side of~\eqref{eq:complement}
has odd valuation. It cannot be a rational square.
\end{proof}

The theorem only uses the residue configuration forced by a negative
odd detector pole. It permits poles of the multiplier on the other
two residue configurations.

\section{A concrete adaptive denominator and its contact order}

Take
\begin{equation}\label{eq:concrete}
              P(v,j)=1,\qquad Q(v,j)=1-v+j,
              \qquad m^*=\frac{1}{1-V^*+J^*}.
\end{equation}
If $u_+,u_-$ are the two slots of $(W^*)^{-1}$, then
\begin{equation}\label{eq:normdenominator}
             1-V^*+J^*=(1-u_+)(1-u_-)
                         =\Nm(1-(W^*)^{-1}).
\end{equation}
At each prepared real character the original slots of $W$ are at
least one. Since $w>0$, the slots of $W^*=(1+w)W$ exceed one.
Thus $0<u_+,u_-<1$, and~\eqref{eq:normdenominator} is strictly
positive. Both it and $m^*$ are whole-algebra units. Also
$Q(0,0)=1$, as required by Theorem~\ref{thm:adaptive}. At the
selected prepared rational character one has $m^*>1$, so
$17(m^*)^2>17>b$.

The exact prescribed first coordinate is
\begin{equation}\label{eq:actualu}
             m^*V^*=\frac{V^*}{1-V^*+J^*}
                         =\frac{\Tr(W^*)}{\Nm(W^*-1)}.
\end{equation}

\begin{proposition}\label{prop:contact}
Let $p\notin S$ be split in $K$, and suppose $v_p(w)=1$.
Write $W_+,W_-$ for its two local slots. Suppose
\[
                         v(W_+)=0,\qquad v(W_-)<0.
\]
At a guarded specialization,
\[
             v_p(b)=1,\qquad v_p(m^*)\le-1.
\]
The tests $h_{k,+}$ on the unit branch are integral. If
\begin{equation}\label{eq:contactnoncancellation}
                     1+\sum_{k=1}^{r}h_{k,+}^2
                              \not\equiv0\pmod p,
\end{equation}
then $v_p(m^*)=-1$. More generally
$v_p(m^*)=-v(W_+^*-1)$.
\end{proposition}
\begin{proof}
Each original factor has nonpositive valuation. Since their product
$W_+$ has valuation zero, every factor on that branch has valuation
zero. The odd weight valuation then gives $v(h_{k,+})\ge0$,
and $W_+\equiv1\pmod p$. Hence
$W_+^*=(1+w)W_+\equiv1\pmod p$. The other slot of
$(W^*)^{-1}$ has positive valuation. Therefore $V^*$ is a unit,
and $b=4w^*/(V^*)^2$ has valuation one.

The factor $1-u_-$ in~\eqref{eq:normdenominator} is a unit,
whereas $1-u_+$ has valuation $v(W_+^*-1)\ge1$.
This proves the general formula and inequality. Expanding the
product modulo $p^2$ gives
\[
               \frac{W_+^*-1}{w}
                   \equiv1+\sum_{k=1}^{r}h_{k,+}^2\pmod p.
\]
Condition~\eqref{eq:contactnoncancellation} makes the last value
a unit, proving the exact order one. Without that condition,
cancellation can increase the order and no equality is asserted.
\end{proof}

\section{Varying admissible norm classes}\label{sec:classes}

\begin{theorem}\label{thm:classes}
Let $\omega$ denote either $w$ or $w^*$, and let $R$ be a fixed
native scalar rational function, taking values in $\mathcal B^\times$
at every prepared point-algebra specialization. Fix a positive rational $S$-unit
$\kappa>1/17$, and put
\begin{equation}\label{eq:varyingclass}
 d_{R,\kappa}
     =\left(R+\frac{\kappa\omega}{R}\right)^2
                                      -\frac{4\omega}{17}.
\end{equation}
At an actual rational character, $v_p(d_{R,\kappa})$ is even
whenever $p\notin S$ and $v_p(\omega)$ is odd.
At every prepared real character, $d_{R,\kappa}>0$, so it is
a unit of the whole point algebra. In particular these classes
can be used with Proposition~\ref{prop:tracelevels} or
Theorem~\ref{thm:adaptive}, respectively, without a new retained
coordinate.
\end{theorem}
\begin{proof}
Write $n=v_p(\omega)$ and $r_0=v_p(R)$. Since $\kappa$ is an
$S$-unit, the two summands inside the square in~\eqref{eq:varyingclass}
have valuations $r_0$ and $n-r_0$. They are distinct because $n$
is odd. The valuation of their squared sum is
$2\min(r_0,n-r_0)<n$. It is strictly below the valuation of
$4\omega/17$, and therefore equals the valuation of
$d_{R,\kappa}$. It is even.

Multiplication by $R^2>0$ at a real character gives
\[
 R^2d_{R,\kappa}
    =\kappa^2\omega^2
       +(2\kappa-4/17)R^2\omega+R^4.
\]
As a quadratic polynomial in the real value of $\omega$, its
leading coefficient is positive and its discriminant is
$16(1-17\kappa)R^4/289<0$. It is strictly positive.
All geometric characters are real before the new root adjunctions,
which proves the whole-algebra unit assertion. The only displayed
denominator is the already guarded unit $R$.
\end{proof}

The choice $\kappa=2/17$ gives the simpler formula
\begin{equation}\label{eq:simplerclass}
                       d_R=R^2+\frac{4\omega^2}{289R^2}.
\end{equation}
Every class~\eqref{eq:varyingclass} supplies the bare norm point
\[
        u_0=R+\kappa\omega/R,\qquad y_0=1,
        \qquad u_0^2-d_{R,\kappa}y_0^2=4\omega/17.
\]
This identity does not prescribe the native first coordinate.

\section{A normalized class with a fixed norm point}

\begin{proposition}\label{prop:normalizedclass}
For the pole-preserving normalization, put
\begin{equation}\label{eq:normalizedclass}
                         d_0=1-\frac{4w^*}{17}.
\end{equation}
This value is a local unit at every $p\notin S$ with $v_p(w^*)$
odd. At every prepared point-algebra character,
\[
                    \frac{16}{17}\le d_0<1.
\]
It is consequently a whole-algebra unit. Its bare norm conic has
the fixed point $(u,y)=(1,1)$.
\end{proposition}
\begin{proof}
Every odd valuation of $w^*$ is positive, so $d_0\equiv1\pmod p$.
For prepared $w>0$, the identity $(1+w)^2-4w=(1-w)^2$ gives
$0<w^*\le1/4$. This proves the real bounds and nonvanishing.
Finally $1-d_0=4w^*/17$ proves the norm-point assertion.
\end{proof}

There is also a class family that can depend on the auxiliary
point data. Let $D,E\in\Z[S^{-1}][s,v,j]$, and assume
$D(0,0,0)$ and $E(0,0,0)$ are $S$-units. Define
\begin{equation}\label{eq:residueclass}
          d=\frac{D(w^*,V^*,J^*)}{E(w^*,V^*,J^*)},
\end{equation}
retaining both numerator and denominator guards.

\begin{proposition}\label{prop:residueclass}
At every $p\notin S$ for which $v_p(b)$ is negative and odd,
$d$ in~\eqref{eq:residueclass} is a local unit. It can therefore
be used in Theorem~\ref{thm:adaptive}. If its numerator and
denominator are positive at every prepared real character, both
are whole-algebra units and their quotient is positive there.
\end{proposition}
\begin{proof}
The proof of Theorem~\ref{thm:adaptive} gives
$(w^*,V^*,J^*)\equiv(0,0,0)\pmod p$. The specialized
polynomials reduce to their nonzero constant coefficients, so both
are local units. The assumed characterwise positivity proves the
last assertion.
\end{proof}

The unit-constant condition is only a local criterion. The
whole-character nonvanishing in the last assertion must still be
proved for the chosen expressions. For example $d_0$ satisfies
both requirements automatically. Both the finite list of such
expressions and all their coefficients are fixed before the input;
no new class-selection parameter is introduced.

\section{All-character complement guards}

\begin{lemma}\label{lem:complementguard}
Let $\mathcal B$ be the native point algebra at a prepared actual
specialization. Suppose $\omega\in\Q^\times$ has even valuation
at $17$, and $m,T\in\mathcal B^\times$. Then
\[
                       17m^2-\omega T^2\in\mathcal B^\times.
\]
\end{lemma}
\begin{proof}
If a geometric character annihilated this element, then
$17/\omega=(\chi(T)/\chi(m))^2\in K^{\times2}$.
A nonzero rational number that is a square in $K$ is a rational
square or twice a rational square: the irrational coefficient of
$(a+b\sqrt2)^2$ vanishes only when $ab=0$. Both classes have
even $17$-adic valuation, whereas $17/\omega$ has odd valuation.
This contradiction proves nonvanishing at every character.
\end{proof}

For the original detector take $T=2/V$ and $\omega=w$.
For the re-expressed detector take $T=2/V^*$ and $\omega=w^*$;
Lemma~\ref{lem:exact} gives
$v_{17}(w^*)\equiv v_{17}(w)\equiv0\pmod2$.
Thus $a-b$ is a whole-algebra unit for all the levels considered,
provided the numerator and denominator units of $m$ are established.
Proposition~\ref{prop:tracelevels} and~\eqref{eq:concrete} supply
two explicit sufficient conditions. For other $P,Q$, that
nonvanishing requirement must be checked; $Q(0,0)$ alone is only
the local soundness condition.

Theorem~\ref{thm:classes}, Proposition~\ref{prop:normalizedclass},
or the stated guards in Proposition~\ref{prop:residueclass} make $d$
a unit. Therefore
$d a(a-b)$ is a whole-algebra unit, and each new quadratic root
adjunction is etale. The norms of every numerator, denominator,
$1+w$, $R$, and $a-b$ that is needed are included in the scalar
guard. Their nonvanishing has been proved on prepared tuples;
none is omitted from the encoding.

\section{Actual rational realization and the variable count}

At a rational character, the complement equation for $a=17m^2$
is equivalent to the prescribed-trace norm equation
\begin{equation}\label{eq:actualnorm}
                  (m\mathcal V)^2-dy^2=4\omega/17,
\end{equation}
where $(\omega,\mathcal V)=(w,V)$ or $(w^*,V^*)$ as appropriate.
The rational conversions are
\[
                  y=\frac{q\mathcal V}{17dm},\qquad
                  q=\frac{17dm y}{\mathcal V}.
\]
All displayed divisors have already been guarded.

Fix one coefficient member of a native source in its sufficient
soundness range. Choose the level and class functions from the preceding families
in advance, with their fixed coefficients, allowing a fixed finite
list of choices. Assume the all-character guards specified above have been
proved for these choices.

\begin{definition}\label{def:realization}
\emph{Simultaneous actual realization} means that for every integer
input in the native completeness range there is an actual prepared
guarded point-stage tuple and an allowed fixed choice of the functions
such that both equations~\eqref{eq:actualnorm} have rational solutions.
The two physical arguments use the same native comparison slope and
the source's prescribed root-dependent height scale. Their traces
are the native traces at that tuple. For a conditional source the
range consists of nonzero integers; zero is supplied by its final
input factor.
\end{definition}

This is a global rational existence assertion. The freely specified
points in Section~\ref{sec:classes}, finite-place solubility, and the guard lemmas do
not establish Definition~\ref{def:realization}.

\begin{theorem}\label{thm:count}
Under the retained source premises and
Definition~\ref{def:realization}, the two pole stages can be replaced
by finite quadratic choices with no retained pole coordinate. The
resulting integer polynomial defines $\Z$ over $\Q$ with the following
total variable counts, including the input:
\begin{center}
\begin{tabular}{lrrrr}\toprule
Source & Point pairs & Input and point coordinates & Slope and encoder & Total\\\midrule
General &9&19&2&21\\
Either conditional &3&7&2&9\\\bottomrule
\end{tabular}
\end{center}
Without Definition~\ref{def:realization}, these counts are not
established defining bounds.
\end{theorem}
\begin{proof}
Retain the input, the $2p$ rational square-abscissa coordinates,
and the single comparison slope, with $p=9$ or $p=3$.
The native point algebra is free of rank $4^p$ over their integer
polynomial ring. Clear the already specified guarded denominators
of each expression $d_i a_i(a_i-b_i)$ and write it as $C_i/D_i$.
Adjoin the two monic roots
\[
                         q_i^2-C_iD_i=0\qquad(i=0,1).
\]
At a rational character, $q_i/D_i$ is precisely a square root of
the corresponding complement. The new finite algebra is free of
rank $4^{p+1}$. Include the inherited point-stage guards and the
norms of all new denominator and unit factors in the scalar guard.
The preceding results make both radicands units at every prepared
character, so the algebra is etale there.

The guarded separator and determinant construction
of~\cite[Section 6]{Unit} uses one encoder coordinate. Its rational
zeros are exactly the rational characters satisfying these guards;
when a scalar guard vanishes its value is one. Polynomial envelopes
and multiplication matrices add no rational parameters.

For soundness, each rational character supplies the unchanged native
detector outputs and the two complement equations. Their local
conclusions exclude negative odd valuations outside the original
exceptional set. The remaining source local extraction and height
arguments are unchanged, since Lemma~\ref{lem:exact} preserves $b$
exactly. The rational input is therefore integral.

For completeness, Definition~\ref{def:realization} supplies the
actual tuple and rational roots. The all-character unit results
verify every new guard and etale root adjunction, so the separator
supplies the encoder. A conditional source includes zero by the same
final factor $X$ as its native construction. A fixed finite union
of choices is encoded by a product of the corresponding determinant
polynomials. There are $1+2p+1+1=2p+3$ variables in each case.
\end{proof}

In particular, one concrete arithmetic target is
\[
 \left(\frac{V_i^*}{1-V_i^*+J_i^*}\right)^2
       -d_{R_i,\kappa_i}y_i^2=4w_i^*/17,
                         \qquad i=0,1,
\]
on actual native tuples at the two coupled arguments. The simpler
class $d_i=1-4w_i^*/17$ is also available. Their simultaneous rational
realization is unproved. The established
general defining-variable counts remain twenty-five and thirteen.

\section{Exact algebraic checks}

The companion verifier
\begin{center}\small
\texttt{computation/pole\_parameters/adaptive\_levels\_v50/check.py}
\end{center}
checks the universal factor, trace, norm, pole-input and norm-equation
identities by exact symbolic arithmetic. A separate implementation of
$\Q(\sqrt2)$ with rational pairs checks native products at split and
inert primes, both signs of odd original weight valuations, the residue
configurations, and the varying-class parity formula. These finite
examples are products of the stated native form; they are not asserted
to arise from elliptic point tuples.

An exact example distinguishes the adaptive multiplier from a locally
integral level at every odd contact. Set
\[
 w=\frac7{16},\qquad h=\frac{4\sqrt2-12}{7},\qquad
 W=1+wh^2=\frac{18-6\sqrt2}{7}.
\]
Then $V=1$, $J=7/36$, $b=7/4$, and
\[
 V^*=\frac{16}{23},\qquad J^*=\frac{448}{4761},\qquad
 1-V^*+J^*=\frac{1897}{4761},\qquad
 m^*=\frac{4761}{1897}.
\]
Since $1897=7\cdot271$ and $7\nmid4761$, the multiplier has
valuation $-1$ at the positive odd contact $7$. This example does not
satisfy the full elliptic preparation hypotheses and is not used as a
completeness witness. The proofs above, rather than the finite checks,
establish the universal local and whole-algebra assertions.

\begin{thebibliography}{9}
\bibitem{Unit}
\emph{A single norm determinant for the 25- and 13-variable constructions},
repository source,
\nolinkurl{papers/finite_algebra/unit_norm_20261114_v28.tex}.
\bibitem{Terminal}
\emph{Smaller norm polynomials with a possibly singular final quadratic extension},
repository source,
\nolinkurl{papers/finite_algebra/terminal_quadratic_root_20261115_v40.tex}.
\bibitem{Complement}
\emph{Unit guards and a finite square-complement reduction for elliptic detectors},
repository source,
\nolinkurl{papers/pole_parameters/guarded_complement_reduction_20261116_v45.tex}.
\bibitem{Finite}
\emph{Finite square-class conditions for native detector inputs},
repository source,
\nolinkurl{papers/pole_parameters/finite_complement_classes_20261116_v46.tex}.
\bibitem{Bounded}
\emph{Bounded complement levels from native detector norms},
repository source,
\nolinkurl{papers/pole_parameters/bounded_norm_levels_20261116_v48.tex}.
\end{thebibliography}
\end{document}
